\section{Discussion}
\label{sec:discussion}

\subsection{Practical implications for solver design}

The abstraction-error results suggest that card-abstraction granularity
above $K\!\approx\!50$ yields diminishing quality improvements on dry
and semi-wet boards within a $1{,}000$-iteration budget, while
action-abstraction error decreases monotonically in $|S|$ with the
single largest drop coming from adding a pot-sized bet ($\{0.25,0.5,0.75\}\!\to\!\{0.25,0.5,0.75,1.0\}$
roughly halves $\epsilon_{\mathrm{abs}}$, from $0.91$ to $0.42$\,bb).
Practitioners building HUNL solvers should prioritize the pot-sized bet
when budget for additional action-tree branches becomes available, and
should expect the practical iteration count needed to enter CFR$^+$'s
asymptotic regime to grow super-linearly with $K$ at deep stacks.

\subsection{Relation to broader game-solving}

The CFR--FTRL connection has implications beyond poker. Any finite
two-player zero-sum extensive-form game can in principle be solved using
the same framework. The abstraction methodology (equity bucketing, action
discretization) generalizes to any game with continuous action spaces or
large state representations.

\subsection{Limitations}

\begin{itemize}
  \item Flop-subgame solving only; preflop and river results may differ.
  \item Abstraction-error measurement uses a proxy ``true'' game, not the
        actual full game.
  \item Single-threaded implementation; parallel CFR may exhibit different
        convergence behavior.
\end{itemize}
